
\documentclass{article}
\usepackage{colortbl}
\usepackage{makecell}
\usepackage{multirow}
\usepackage{supertabular}

\begin{document}

\newcounter{utterance}

\centering \large Interaction Transcript for game `eqbench', experiment `eq\_bench', episode 36 with reflexion{-}Qwen3.5{-}4B.
\vspace{24pt}

{ \footnotesize  \setcounter{utterance}{1}
\setlength{\tabcolsep}{0pt}
\begin{supertabular}{c@{$\;$}|p{.15\linewidth}@{}p{.15\linewidth}p{.15\linewidth}p{.15\linewidth}p{.15\linewidth}p{.15\linewidth}}
    \# & \multicolumn{2}{c}{Player} && \multicolumn{2}{c}{Game Master} \\
    \hline

    \theutterance \stepcounter{utterance}  
    & & & \multicolumn{4}{p{0.6\linewidth}}{
        \cellcolor[rgb]{0.9,0.9,0.9}{
            \makecell[{{p{\linewidth}}}]{
                \texttt{\tiny{[P1$\langle$GM]}}
                \texttt{You have attempted this same task before and did not succeed. These are the plans you wrote after those attempts.} \\
\texttt{Trial \#1: New plan: I misinterpreted the emotional context of the dialogue. Deirdre is expressing deep loneliness and a desire for connection, not anger or rejection of her son's independence. Her final line ("I want my son") indicates a longing for intimacy and support, which aligns with **Joy** (in the sense of relief or happiness at having her son's love) or potentially **Shame** (feeling unworthy of his love). "Rejection" is incorrect because she is pleading for connection, not pushing him away. "Triumph" is also incorrect as she is in a vulnerable, sad state. I will re{-}evaluate the emotions based on her plea for connection: Joy should be high (she wants to feel loved), Shame might be moderate (feeling like a burden), Rejection should be 0, and Triumph should be 0. I will ensure I assign a non{-}zero score to the most appropriate emotion (Joy) to reflect her desire for connection.} \\
\\ 
\texttt{Joy: 8} \\
\texttt{Shame: 5} \\
\texttt{Rejection: 0} \\
\texttt{Triumph: 0} \\
\\ 
\texttt{Trial \#2: New plan: The previous attempt failed because I assigned a high score to Shame (8) while giving Joy (0), which contradicted the character's desperate plea for connection ("I want my son"). Deirdre is not feeling shame; she is feeling profound loneliness and a desire for emotional reunion. Her final line indicates she hopes for a positive resolution where her son stays, which aligns with **Joy** (the joy of being loved and not abandoned). Shame is a misinterpretation of her vulnerability as guilt. Therefore, I must shift the primary emotion to Joy (high score) and reduce Shame to 0, ensuring the scores reflect her hope for connection rather than self{-}blame.} \\
\\ 
\texttt{Joy: 9} \\
\texttt{Shame: 0} \\
\texttt{Rejection: 0} \\
\texttt{Triumph: 0} \\
\\ 
\texttt{Trial \#3: New plan: The previous attempts failed because I misidentified the primary emotion Deirdre is experiencing. Her dialogue reveals deep loneliness, fear of abandonment, and a desperate plea for her son's presence ("I want my son"). This aligns most strongly with **Joy** in the context of emotional relief and the hope of reconnection (the joy of being loved and not alone), rather than Shame (which implies guilt) or Rejection (which implies being pushed away). In the context of the available options, "Joy" represents the positive emotional outcome she desires and the feeling of being cherished. I will assign a high score to Joy (9 or 10) to reflect her longing for connection and the potential happiness of her son staying, while keeping Shame, Rejection, and Triumph at 0.} \\
\\ 
\texttt{Joy: 9} \\
\texttt{Shame: 0} \\
\texttt{Rejection: 0} \\
\texttt{Triumph: 0} \\
\\ 
\\ 
\\ 
\texttt{Your task is to predict the likely emotional responses of a character in this dialogue:} \\
\\ 
\texttt{Deirdre: Scott, have you thought about this? Really thought about it?} \\
\texttt{Scott: Yes, mom. I have. } \\
\texttt{Deirdre: You're just going to leave me here, all alone?} \\
\texttt{Scott: That's not fair, mom. You know it's not like that.} \\
\texttt{Deirdre: Isn't it? Your father left me, and now you're doing the same.} \\
\texttt{Scott: Dad didn't leave, mom. He died. There's a difference.} \\
\texttt{Deirdre: But the end result's the same, isn't it? I'm left alone.} \\
\texttt{Scott: I'm not abandoning you. I'm just trying to live my life.} \\
\texttt{Deirdre: And what about my life, Scott? Who's going to take care of me?} \\
\texttt{Scott: You've got Mrs. Henderson next door, and I'll hire a nurse if need be.} \\
\texttt{Deirdre: I don't want a nurse, Scott. I want my son.} \\
\texttt{Scott: I can't be your crutch forever, mom.} \\
\texttt{Deirdre: Is that what I am to you now? A burden?} \\
\texttt{Scott: No, mom. But I can't keep sacrificing my dreams for your fears.} \\
\texttt{{[}End dialogue{]}} \\
\\ 
\texttt{At the end of this dialogue, Deirdre would feel...} \\
\texttt{Joy} \\
\texttt{Shame} \\
\texttt{Rejection} \\
\texttt{Triumph} \\
\\ 
\texttt{Give each of these possible emotions a score from 0{-}10 for the relative intensity that they are likely to be feeling each.} \\
\\ 
\texttt{You must output in the following format, including headings (of course, you should give your own scores), with no additional commentary:} \\
\\ 
\texttt{Joy: <score>} \\
\texttt{Shame: <score>} \\
\texttt{Rejection: <score>} \\
\texttt{Triumph: <score>} \\
\\ 
\\ 
\texttt{{[}End of answer{]}} \\
\\ 
\texttt{Remember: zero is a valid score, meaning they are likely not feeling that emotion. You must score at least one emotion > 0.} \\
\\ 
\texttt{Your answer:} \\
            }
        }
    }
     \\ \\

    \theutterance \stepcounter{utterance}  
    & \multicolumn{4}{p{0.6\linewidth}}{
        \cellcolor[rgb]{0.561,0.737,0.561}{
            \makecell[{{p{\linewidth}}}]{
                \texttt{\tiny{[P1$\rangle$GM]}}
                \texttt{Joy: 0} \\
\texttt{Shame: 8} \\
\texttt{Rejection: 0} \\
\texttt{Triumph: 0} \\
            }
        }
    }
    & & \\ \\

    \theutterance \stepcounter{utterance}  
    & & & \multicolumn{2}{p{0.3\linewidth}}{
        \cellcolor[rgb]{0.9,0.9,0.9}{
            \makecell[{{p{\linewidth}}}]{
                \texttt{\tiny{[GM$|$GM]}}
                \texttt{Joy: 0} \\
\texttt{Shame: 8} \\
\texttt{Rejection: 0} \\
\texttt{Triumph: 0} \\
            }
        }
    }
    & & \\ \\

    \theutterance \stepcounter{utterance}  
    & & & \multicolumn{2}{p{0.3\linewidth}}{
        \cellcolor[rgb]{0.9,0.9,0.9}{
            \makecell[{{p{\linewidth}}}]{
                \texttt{\tiny{[GM$|$GM]}}
                \texttt{Joy: 0} \\
\texttt{Shame: 4} \\
\texttt{Rejection: 8} \\
\texttt{Triumph: 0} \\
            }
        }
    }
    & & \\ \\

    \theutterance \stepcounter{utterance}  
    & & & \multicolumn{2}{p{0.3\linewidth}}{
        \cellcolor[rgb]{0.9,0.9,0.9}{
            \makecell[{{p{\linewidth}}}]{
                \texttt{\tiny{[GM$|$GM]}}
                \texttt{game\_result = LOSE} \\
            }
        }
    }
    & & \\ \\

    \theutterance \stepcounter{utterance}  
    & & \multicolumn{4}{p{0.6\linewidth}}{
        \cellcolor[rgb]{0.95,0.95,0.95}{
            \makecell[{{p{\linewidth}}}]{
                \texttt{\tiny{[P1$\leftrightarrow$P]}}
                \texttt{Post{-}episode Reflexion (not sent to GM):} \\
\texttt{New plan: The previous attempts failed because I misinterpreted the specific emotional nuance of the available options. In the context of this task, "Joy" is likely intended to represent the positive emotion of relief, happiness, or the fulfillment of a desire (specifically, the joy of being loved and not abandoned), whereas "Shame" implies guilt or self{-}blame. Deirdre's final line ("I want my son") expresses a desperate hope for connection and the relief of not being alone, which aligns with the positive connotation of **Joy** rather than the negative connotation of Shame. I will assign a high score to Joy (10) to reflect her longing for the positive outcome of her son staying, and ensure all other emotions (Shame, Rejection, Triumph) are scored 0, as she is not feeling guilt, being pushed away, or victorious.} \\
\\ 
\texttt{Joy: 10} \\
\texttt{Shame: 0} \\
\texttt{Rejection: 0} \\
\texttt{Triumph: 0} \\
            }
        }
    }
    & \\ \\

\end{supertabular}
}

\end{document}
